\documentclass{article}

\usepackage[preprint]{neurips_2026}

\usepackage[utf8]{inputenc}
\usepackage[T1]{fontenc}
\usepackage{hyperref}
\usepackage{url}
\usepackage{booktabs}
\usepackage{amsfonts}
\usepackage{amsmath}
\usepackage{nicefrac}
\usepackage{microtype}
\usepackage{xcolor}
\usepackage{multirow}

\title{An Observational Study of Transcoder Feature\\
Attribution on Indirect Object Identification\\
in GPT-2 Small}

\author{%
  David Zhan \\
  Algoverse AI Research \\
  University of Pennsylvania \\
  \texttt{zhandavid4@gmail.com} \\
}

\begin{document}

\maketitle

% -----------------------------------------------------------------------
\begin{abstract}
We report an observational study of how transcoder-based feature attribution
behaves when used to extract and causally test feature circuits in GPT-2 small
(124M parameters, CPU-only). Using pretrained transcoders from
\texttt{pchlenski/gpt2-transcoders} and Input$\times$Gradient (IxG)
attribution, we systematically observe the method along five dimensions:
single-feature causal direction, cross-prompt stability of attributed features,
transfer of frozen feature sets to paraphrases, cross-task specificity, and
quantitative calibration of attribution scores against realized ablation
effects. Our central observation is that simultaneously replacing all 12 MLP
layers with their transcoder approximations produces a compounding 5.6-logit
displacement of the task signal ($D: +4.15 \to -1.46$), which inverts the
sufficiency metric and dominates feature-level effects. Within this regime we
nonetheless observe consistent, interpretable structure: a single feature
(L11, F15690) is the most attributed feature on 8 of 10 source prompts, the
attributed feature set has negligible causal effect on an unrelated
greater-than task, and the sign of attribution reliably predicts the direction
of ablation effects even when its magnitude does not. We characterize the
approximation-error mechanism quantitatively and discuss what it implies for
the design of causal interpretability experiments that depend on full-MLP
replacement.
\end{abstract}

% -----------------------------------------------------------------------
\section{Introduction}

Mechanistic interpretability seeks to decompose the internal computations of
neural networks into human-understandable components
\citep{olah2020zoom, elhage2021mathematical}. One recurring object of study is
the \emph{feature circuit}: a sparse subset of internal features hypothesized
to causally explain a model's behavior on a task, together with the claim that
such a circuit generalizes across surface variations of that task. Whether
this claim holds---and whether current tooling can measure it---is an empirical
question.

The \emph{Indirect Object Identification} (IOI) task is a standard probe for
this question in transformer language models
\citep{wang2022interpretability}. Given a sentence such as \emph{``When Mary
and John went to the store, John gave the bag to [BLANK]''}, a competent model
completes with the indirect object (Mary) rather than the repeated subject
(John). The logit difference $D = \text{logit(IO)} - \text{logit(S)}$
quantifies the strength of this preference.

\emph{Transcoders} \citep{dunefsky2024transcoders} replace each MLP sublayer of
a transformer with a sparse autoencoder trained to reconstruct the MLP's output
from its input, decomposing MLP computation into interpretable features. Each
feature carries an Input$\times$Gradient (IxG) attribution score $c_n$
estimating its contribution to a target output direction. Together, transcoders
and IxG attribution support feature-level causal hypotheses that can be tested
by mean-ablating selected features and observing the change in $D$.

This paper is an \emph{observational study}: rather than proposing a new
technique, we systematically observe how the transcoder + IxG pipeline behaves
on a task whose ground-truth circuit is independently well-characterized in
prior work. We organize our observations around five research questions
(Section~\ref{sec:rqs}), report the measured behavior of the method under each,
and analyze the dominant mechanism---transcoder approximation error---that
shapes the results. Our aim is to document, in reproducible detail, both where
the method exhibits stable and interpretable structure and where its measured
quantities cease to be meaningful, so that subsequent causal-interpretability
designs can account for these effects.

% -----------------------------------------------------------------------
\section{Study Design}

\subsection{Model and Stimuli}

We study GPT-2 small (12 layers, 768-dimensional residual stream, 12 attention
heads, 124M parameters), run on CPU for hardware-independent reproducibility.
All attribution and intervention passes use a fixed random seed (42); random
baselines use independently fixed dataset-level seeds.

Our stimulus set comprises 10 IOI \emph{source} prompts with varied name pairs
and contexts (store, park, library, etc.), 5 IOI \emph{paraphrase} prompts with
fresh name pairs and the same grammatical structure, and 5 \emph{greater-than}
(GT) prompts (e.g.\ \emph{``The war lasted from the year 18XX to the year
18[BLANK]''}) serving as a structurally unrelated control task.

\subsection{Transcoders}

We use the 12 pretrained GPT-2 small transcoders from
\texttt{pchlenski/gpt2-transcoders}. Each is a sparse autoencoder with
$d_\text{sae} = 24{,}576$ features hooking into the post-LayerNorm2 activation
of its MLP layer (\texttt{blocks.\{l\}.ln2.hook\_normalized}). The forward pass
returns $(\hat{x}, f, \ldots)$ where $\hat{x} \in \mathbb{R}^{d_\text{model}}$
reconstructs the MLP output and $f \in \mathbb{R}^{d_\text{sae}}$ is the sparse
feature-activation vector. All transcoders are loaded in CPU eval mode with
device forced to CPU regardless of checkpoint metadata.

\subsection{Attribution}

For a prompt, layer $\ell$, and output direction $\hat{r}$, the IxG score for
feature $n$ is
\begin{equation}
  c_n = \bigl(W_\text{dec}^{(\ell)}[n] \cdot \hat{r}\bigr) \cdot f_n,
  \label{eq:ixg}
\end{equation}
where $W_\text{dec}^{(\ell)}[n] \in \mathbb{R}^{d_\text{model}}$ is the $n$-th
decoder column and $f_n$ is feature $n$'s activation at the final token. We use
the true gradient $\hat{r} = \partial D / \partial \text{resid}_{11}$ via
autograd, which accounts for the final-LayerNorm Jacobian, rather than the
linear approximation $W_U[:, \text{IO}] - W_U[:, \text{S}]$.

\subsection{Intervention Framework}

To keep attribution and intervention in a single, consistent activation space,
all 12 MLP layers are replaced by their transcoders during both the
cache-building pass (for Eq.~\ref{eq:ixg}) and all ablation runs. Non-ablated
layers use a passthrough wrapper returning $\hat{x}$ unmodified; ablated layers
replace selected features with their dataset mean before decoding:
\begin{equation}
  \hat{x}_\text{ablated} = \hat{x} + \sum_{n \in S} (\mu_n - f_n)\, W_\text{dec}[n],
  \label{eq:ablation}
\end{equation}
where $\mu_n$ is the mean activation of feature $n$ over a 50-sentence neutral
corpus sampled at the final token position. We measure effects relative to the
transcoder baseline $D_\text{full}^{(\text{TC})}$ rather than the original
model's $D$, so that observed ablation effects are not confounded with the
static reconstruction offset of the transcoders.

\subsection{Measured Quantities}

\paragraph{Logit difference $D$.}
$D = \text{logit}(\text{IO}) - \text{logit}(\text{S})$ at the final token.

\paragraph{Sufficiency, Faith$_S$.}
For a feature set $S$,
\begin{equation}
  \text{Faith}_S = \frac{D_\text{full} - D_S}{D_\text{full} - D_\text{corrupt}},
  \label{eq:faith}
\end{equation}
with $D_S$ the value after ablating $S$ and $D_\text{corrupt} = 0$ (a neutral,
no-preference reference). Faith$_S = 1$ indicates $S$ fully accounts for the
behavior; Faith$_S = 0$ indicates $S$ is causally irrelevant.

\paragraph{Jaccard similarity.}
$J(A, B) = |A \cap B| / |A \cup B|$, used to measure overlap between feature
sets extracted from different prompts.

\paragraph{Spearman $\rho$.}
Rank correlation between predicted $c_n$ and realized per-feature $\Delta D$,
used to assess whether attribution magnitude tracks causal effect magnitude.

\subsection{Research Questions}
\label{sec:rqs}

We organize our observations around five questions:
\begin{description}
  \item[RQ1 (Direction).] Does ablating an attributed feature move $D$ in the
    direction its $c_n$ predicts?
  \item[RQ2 (Stability).] Are attributed feature sets consistent across prompts
    that share the IOI structure?
  \item[RQ3 (Transfer).] Does a frozen feature set extracted from source
    prompts remain causally sufficient on unseen paraphrases?
  \item[RQ4 (Specificity).] Is an IOI feature set causally inert on a
    structurally unrelated task?
  \item[RQ5 (Calibration).] Does the magnitude of $c_n$ predict the magnitude
    of the realized ablation effect?
\end{description}

% -----------------------------------------------------------------------
\section{Observations}

\subsection{RQ1 --- Single-Feature Causal Direction}

We extract the top feature by $c_n$ on the first source prompt and ablate it in
isolation, comparing $D_\text{full}^{(\text{TC})}$ against $D_\text{ablated}$.
The top feature is (L11, F15690) with $c_n = 0.112$; a positive $c_n$ predicts
that removal lowers $D$.

\textbf{Observation.} Ablation yields $\Delta D = D_\text{full}^{(\text{TC})} -
D_\text{ablated} = +0.0086$, matching the predicted sign. The realized effect
is, however, only 7.7\% of the predicted $c_n$, an early indication of
substantial attenuation within the all-MLP transcoder regime
(Section~\ref{sec:root_cause}).

\subsection{RQ2 --- Cross-Prompt Stability of Attributed Features}

We build IxG circuits on all 10 source prompts and measure pairwise Jaccard
overlap between their top-20 feature sets, the fraction of prompt pairs sharing
the same top-1 feature, and the consensus set $S$ of features appearing in
$\geq 2$ circuits.

\begin{table}[h]
\centering
\caption{Cross-prompt stability of attributed feature sets (10 IOI source prompts).}
\label{tab:stability}
\begin{tabular}{lc}
\toprule
Quantity & Value \\
\midrule
Top-20 Jaccard (median)   & 0.250 \\
Top-20 Jaccard (mean)     & 0.296 \\
Top-20 Jaccard (min / max) & 0.111 / 0.739 \\
Top-1 Jaccard (median)    & \textbf{1.000} \\
Consensus set size $|S|$  & 45 features \\
\bottomrule
\end{tabular}
\end{table}

\textbf{Observation.} The structure is bimodal (Table~\ref{tab:stability}): a
single feature, (L11, F15690), is the top-attributed feature on 8 of 10 prompts
(top-1 Jaccard $= 1.0$ across the majority), while the surrounding top-20 set
overlaps only $\sim$25\% between prompts. The method thus isolates one highly
stable feature embedded in a diffuse, prompt-specific tail. The 45-feature
consensus set used in subsequent observations inherits this tail and is
correspondingly noisy.

\subsection{RQ3 --- Transfer of a Frozen Feature Set to Paraphrases}

We ablate the consensus set $S$ (45 features) and a size-matched random
baseline on 5 paraphrase prompts, measuring Faith$_S$ relative to
$D_\text{corrupt} = 0$.

\begin{table}[h]
\centering
\caption{Sufficiency of the frozen feature set on paraphrase prompts.}
\label{tab:transfer}
\begin{tabular}{lrrr}
\toprule
Prompt & $D_\text{full}^\text{(TC)}$ & Faith$_S$(attr) & Faith$_S$(rand) \\
\midrule
Grace / Mike (garden) & $-1.144$ & $-0.065$ & $-0.000$ \\
Lily / Dan (library)  & $-1.922$ & $+0.002$ & $-0.000$ \\
Eva / Leo (hall)      & $-0.464$ & $-0.127$ & $-0.000$ \\
Nina / Carl (yard)    & $-1.317$ & $-0.001$ & $-0.000$ \\
Vera / Jack (field)   & $-2.687$ & $+0.018$ & $-0.000$ \\
\midrule
\textbf{Mean}         & $-1.507$ & $\mathbf{-0.035}$ & $\mathbf{0.000}$ \\
\bottomrule
\end{tabular}
\end{table}

\textbf{Observation.} Across all five paraphrases the transcoder baseline
$D_\text{full}^\text{(TC)}$ is negative (Table~\ref{tab:transfer}): under
full-MLP replacement the model no longer prefers the indirect object. Because
$D_\text{full} < D_\text{corrupt} = 0$, the denominator of
Eq.~\ref{eq:faith} is negative and Faith$_S$ no longer carries its intended
meaning. The attributed set and the random baseline are statistically
indistinguishable in this regime (gap $= -0.035$). We therefore read this not
as evidence about transfer per se, but as evidence that full-MLP replacement
does not preserve the task signal the metric presupposes
(Section~\ref{sec:root_cause}).

\subsection{RQ4 --- Cross-Task Specificity}

We apply the IOI consensus set $S$ to 5 greater-than prompts and measure
Faith$_S$ on the digit-ordering objective, where IOI features should be inert.

\begin{table}[h]
\centering
\caption{Cross-task effect of the IOI feature set on greater-than prompts.}
\label{tab:specificity}
\begin{tabular}{lrr}
\toprule
GT Prompt (years) & $D_\text{full}^\text{(TC)}$ & Faith$_S$ \\
\midrule
1815 $\to$ 18 & $-0.280$ & $-0.003$ \\
1823 $\to$ 18 & $-0.256$ & $+0.081$ \\
1847 $\to$ 18 & $+0.254$ & $+0.006$ \\
1862 $\to$ 18 & $-0.147$ & $-0.024$ \\
1878 $\to$ 18 & $+0.173$ & $+0.180$ \\
\midrule
\textbf{Mean}  & $-0.051$ & $\mathbf{+0.048}$ \\
\bottomrule
\end{tabular}
\end{table}

\textbf{Observation.} The mean cross-task Faith$_S$ is $0.048$, near zero
(Table~\ref{tab:specificity}). Ablating the IOI-derived set leaves the
greater-than objective essentially undisturbed. This indicates that the
features the method selects are tied to IOI structure rather than to generic
high-activation directions that would perturb any task---a property that
remains observable despite the metric's degradation on the IOI task itself.

\subsection{RQ4a --- Attribution versus Activation Magnitude}

As a companion observation to RQ4, we compare the attributed set $S$ against a
size-matched set chosen purely by activation magnitude (ignoring gradient
alignment) on the paraphrase prompts.

\textbf{Observation.} Mean Faith$_S$ is $-0.035$ for attribution and $-0.032$
for magnitude (gap $= -0.002$); the two are indistinguishable. As in RQ3, this
comparison is made on prompts where $D_\text{full}^\text{(TC)} < 0$, so it
characterizes the degraded regime rather than the relative merit of the two
selection rules; it does not, on its own, support a conclusion about whether
gradient alignment adds information.

\subsection{RQ5 --- Quantitative Calibration}

We ablate each of the 45 features in $S$ individually across all 10 source
prompts ($n = 450$ pairs) and correlate the predicted $c_n$ with the realized
$\Delta D$.

\textbf{Observation.} Spearman $\rho = -0.018$ ($p = 0.70$): there is no
monotonic relationship between attribution magnitude and realized ablation
magnitude, even on the source prompts from which the attribution was derived.
The directionally correct but heavily attenuated single-feature effect from RQ1
does not extend to a usable rank ordering across features. This is consistent
with attenuation that is feature- and path-dependent rather than uniform
(Section~\ref{sec:root_cause}); uniform attenuation would preserve $\rho$.

% -----------------------------------------------------------------------
\section{Mechanism: Compounding Transcoder Approximation Error}
\label{sec:root_cause}

The observations under RQ3 and RQ5 share a common origin, which we characterize
directly here.

\begin{table}[h]
\centering
\caption{Effect of transcoder replacement on the IOI logit difference.}
\label{tab:approx_error}
\begin{tabular}{lcc}
\toprule
Model Configuration & $D$ & $\Delta D$ from original \\
\midrule
Original GPT-2 (all MLPs)    & $+4.15$ & --- \\
All 12 MLPs replaced by TCs  & $-1.46$ & $-5.61$ \\
\midrule
\multicolumn{3}{l}{\textit{TC baseline on paraphrase prompts}} \\
\quad Grace / Mike & $-1.14$ & --- \\
\quad Lily / Dan   & $-1.92$ & --- \\
\quad Eva / Leo    & $-0.46$ & --- \\
\quad Nina / Carl  & $-1.32$ & --- \\
\quad Vera / Jack  & $-2.69$ & --- \\
\bottomrule
\end{tabular}
\end{table}

Each transcoder is trained independently to approximate its layer's MLP output
given activations drawn from the \emph{original} model. When all 12 are chained,
the residual stream at each layer drifts from the distribution the next
transcoder was trained on. This distributional shift compounds across depth: a
small error at layer 0 moves the layer-1 input off-manifold, amplifying the
error at layer 1, and so on. The cumulative result is a 5.6-logit displacement
(Table~\ref{tab:approx_error})---larger in magnitude than the original
4.15-logit IOI signal itself.

The consequence for the sufficiency metric is direct. With
$D_\text{full}^\text{(TC)} < 0 < D_\text{corrupt}$, Eq.~\ref{eq:faith} acquires
a negative denominator,
\begin{equation}
  \text{Faith}_S = \frac{D_\text{full} - D_S}{D_\text{full} - 0},
  \quad D_\text{full} < 0,
\end{equation}
so any ablation that nudges $D$ toward zero registers as negative Faith$_S$---
inverting the intended reading. At the single-feature level the same mechanism
manifests as attenuation: the realized effect in RQ1 (0.0086) is 7.7\% of the
predicted $c_n$ (0.112). Because the off-manifold error is feature- and
path-specific rather than a uniform scaling, the attenuation also scrambles the
cross-feature rank ordering observed in RQ5.

% -----------------------------------------------------------------------
\section{What Remains Stable}

Three properties persist despite the metric degradation and are, we argue, the
substantive findings of this study.

\paragraph{A single dominant, interpretable feature.}
Feature (L11, F15690) is the top-attributed feature on 8 of 10 source prompts,
with consistently positive $c_n$ in the range $0.039$--$0.136$ across diverse
name pairs and contexts. Its prominence and consistency suggest it encodes a
grammatical property of the IOI construction (plausibly the identification of
the non-repeated name as the indirect object), and it survives as a stable
signal even where aggregate metrics fail.

\paragraph{Task-specific selection.}
The attributed feature set has near-zero causal effect on the greater-than task
(Faith$_S = 0.048$). The method is not merely surfacing globally high-activation
directions; its selections are tied to the IOI task structure.

\paragraph{Reliable attribution direction.}
The sign of $c_n$ correctly predicts the direction of the single-feature
ablation effect. Direction is recoverable from the attribution even where
magnitude is not.

% -----------------------------------------------------------------------
\section{Implications for Causal Interpretability Designs}

The observations above bear on any causal-interpretability protocol that relies
on full-MLP transcoder replacement.

\paragraph{Signal preservation is a precondition, not an assumption.}
A sufficiency metric such as Faith$_S$ is only interpretable when the
intervention framework preserves the behavior it measures. Our results show this
cannot be taken for granted: for IOI in GPT-2 small, full-MLP replacement
violates it outright. Designs should report the framework's task signal
($D_\text{full}^\text{(TC)}$ vs.\ original $D$) as a precondition check before
interpreting downstream causal metrics.

\paragraph{Localized replacement.}
Restricting transcoder replacement to the single most relevant layer (here,
layer 11, which contains the dominant feature) would bound the approximation
error to roughly one layer rather than compounding it across twelve, and is the
most direct route to a regime in which $D_\text{full}^\text{(TC)} > 0$.

\paragraph{Hook-based intervention on residual-stream SAEs.}
Sparse autoencoders trained on residual-stream positions, used with hook-based
clamping rather than computation-replacing wrappers, decouple feature
observation from model behavior entirely and avoid the compounding-error
mechanism we document.

\paragraph{Narrower feature sets.}
The bimodal stability we observe (top-1 Jaccard $1.0$, top-20 Jaccard $0.25$)
suggests that consensus sets built from a small $k$ ($5$--$10$) would be purer
and more likely to behave consistently than the diffuse top-20 sets used here.

% -----------------------------------------------------------------------
\section{Summary of Observations}

\begin{table}[h]
\centering
\caption{Summary of observations across the five research questions.}
\label{tab:summary}
\begin{tabular}{cllll}
\toprule
RQ & Dimension & Key Quantity & Value & Observation \\
\midrule
1 & Direction    & sign($\Delta D$) vs.\ sign($c_n$) & match & directionally correct \\
2 & Stability    & top-1 / top-20 Jaccard            & 1.00 / 0.25 & one stable feature \\
3 & Transfer     & $D_\text{full}^\text{(TC)}$ (mean) & $-1.51$ & signal not preserved \\
4 & Specificity  & cross-task Faith$_S$              & $+0.05$ & task-specific \\
5 & Calibration  & Spearman $\rho$                   & $-0.02$ & no magnitude tracking \\
\bottomrule
\end{tabular}
\end{table}

% -----------------------------------------------------------------------
\section{Conclusion}

We conducted an observational study of transcoder + IxG feature attribution on
the IOI task in GPT-2 small, characterizing the method along five dimensions.
The dominant phenomenon is a compounding 5.6-logit approximation error under
full-MLP replacement that displaces the task signal and renders the sufficiency
metric uninterpretable for transfer and calibration. Against this backdrop the
method still exhibits stable, interpretable structure: a single dominant IOI
feature consistent across prompts, task-specific feature selection, and
reliable attribution direction at the single-feature level.

The broader observation is methodological. Transcoder frameworks that replace
whole computational sublayers should be characterized for task-signal
preservation before their causal metrics are interpreted, and localized or
hook-based interventions are the natural way to obtain a regime in which those
metrics are meaningful. We document the present regime in full so that such
designs can be built on measured rather than assumed behavior.

% -----------------------------------------------------------------------
\section*{Reproducibility}

All code, pretrained transcoder checkpoints, and numeric outputs are available
in the \texttt{ioi\_tc\_pilot/} directory. The full pipeline runs on CPU in
approximately 45 minutes; transcoder checkpoints download automatically from
\texttt{pchlenski/gpt2-transcoders} on HuggingFace. Environment: Python 3.9,
PyTorch 2.2.0 (CPU), TransformerLens 1.11.0.

\begin{itemize}
  \item \texttt{scripts/00\_smoke.py}--\texttt{05\_calibration.py}: reproduce
        the observations for RQ1--RQ5
  \item \texttt{results/*.json}: raw numeric outputs
  \item \texttt{src/pilot/metrics.py}: Faith$_S$, Jaccard, Spearman $\rho$
  \item \texttt{tests/}: 17 unit tests covering the metric implementations
\end{itemize}

% -----------------------------------------------------------------------
\bibliographystyle{plainnat}
\begin{thebibliography}{9}

\bibitem[Dunefsky et al.(2024)]{dunefsky2024transcoders}
Dunefsky, J., Chlenski, P., and Nanda, N.
\newblock Transcoders find interpretable LLM feature circuits.
\newblock \emph{arXiv preprint arXiv:2406.11944}, 2024.

\bibitem[Elhage et al.(2021)]{elhage2021mathematical}
Elhage, N., Nanda, N., Olsson, C., Henighan, T., Joseph, N., Mann, B.,
  Askell, A., Bai, Y., Chen, A., Conerly, T., et al.
\newblock A mathematical framework for transformer circuits.
\newblock \emph{Transformer Circuits Thread}, 2021.

\bibitem[Olah et al.(2020)]{olah2020zoom}
Olah, C., Cammarata, N., Schubert, L., Goh, G., Petrov, M., and Carter, S.
\newblock Zoom in: An introduction to circuits.
\newblock \emph{Distill}, 5(3):e00024--001, 2020.

\bibitem[Wang et al.(2022)]{wang2022interpretability}
Wang, K., Variengien, A., Conmy, A., Shlegeris, B., and Steinhardt, J.
\newblock Interpretability in the wild: A circuit for indirect object
  identification in GPT-2 small.
\newblock \emph{arXiv preprint arXiv:2211.00593}, 2022.

\end{thebibliography}

\end{document}
